\section{Introduction}
\label{sec:introduction}

Clinical and population-health datasets are difficult to share and model. They combine sensitive attributes, mixed numerical and categorical variables, missing data, skewed distributions, and outcomes that are often rare but clinically decisive. Synthetic tabular data are increasingly proposed as a way to widen access to such data and to support safer collaboration between institutions, but synthesis does not by itself guarantee privacy, fairness, or downstream usefulness \citep{goncalves2020,hernandez2023,nanevski2026}. A generator can reproduce marginal distributions well while weakening the decision boundary for a rare outcome, or it can improve apparent classifier performance by exaggerating relationships that do not generalise to real patients. These failure modes are especially consequential when the positive class represents death, disease, or an underserved population.

We treat synthetic data as a purpose-specific artefact rather than a single-quality object: statistical fidelity, predictive utility, minority-class preservation, causal plausibility, and privacy are related but distinct properties, and a dataset can satisfy one while failing another. This view was not our starting point but the result of a practical failure encountered early in the project. Changing class counts through conventional resampling, CTGAN-based minority augmentation, and increasingly elaborate causal-anchor filtering did not produce a reliable, substantial improvement in prediction. What these attempts had in common is that they added or removed observations without directly addressing the overlap between outcome classes.

That failure reframed the research question. Rather than asking only how many minority records should be generated, we asked whether a generator can be made to preserve the class-conditional structure that makes a rare outcome learnable in the first place. Explainability and nearest-neighbour analyses supported this reframing: they pointed to weak class separability, not imbalance alone, as a plausible limiting mechanism. This diagnosis motivated \sepa{}, a post-generation selection method that draws an oversized CTGAN candidate pool and explicitly selects records according to realism, label separability, and domain-anchor compatibility, while preserving the real class counts.

\subsection{Research questions}

The overarching question is whether class-conditional separability can be controlled during synthetic tabular health-data generation so as to improve performance on unseen real patients, without producing unrealistic, clinically implausible, or privacy-sensitive synthetic populations. Concretely:
\begin{itemize}[leftmargin=*]
\item Why do class rebalancing and unconstrained synthetic augmentation yield limited or inconsistent gains when outcome classes overlap?
\item Does explicit separability-aware selection improve synthetic-to-real (\tstr{}) utility, and is the effect stable across datasets?
\item When such a gain occurs, is it attributable to separability pressure, to domain-anchor pressure, or to their interaction?
\end{itemize}
The experiments reported below test three directional expectations. First, balancing class counts without changing class-conditional overlap should not by itself yield a large, reliable utility gain. Second, positive separability pressure should increase held-out-real Macro-F1 relative to realism-only selection while leaving global fidelity broadly unchanged. Third, any separability effect should be dataset dependent, with smaller and less stable gains when the positive class is small or heterogeneous. A further, harder question -- whether an observed gain reflects genuine recovery of class-conditional signal or merely the selection of easier synthetic prototypes -- cannot be resolved by the present design and is treated as an open problem in Section~\ref{sec:limitations}.

\subsection{Empirical strategy}

The empirical programme follows the reasoning above rather than a fixed catalogue of models. Early rebalancing and CTGAN-augmentation studies establish that the problem is not solved by frequency alone. A five-fold domain-anchored causal-plausibility filtering (\dacpf{}) experiment tests whether clinically motivated candidate filtering supplies the missing structure. Diagnostic analyses then characterise the class geometry that remains after augmentation. The resulting \sepa{} formulation is evaluated across a fixed-$\lambda$ grid and, more decisively, in a predeclared replicated $2^2$ factorial experiment that isolates separability pressure, anchor pressure, their interaction, and residual error. The factorial study is the paper's central test of mechanism: it asks not merely whether a selected synthetic dataset performs well, but which component of the selection rule is responsible.

In summary, this paper makes the following contributions:
\begin{enumerate}[leftmargin=*]
\item We identify class-conditional learnability, rather than class balance or global fidelity, as a distinct and previously under-specified objective for synthetic tabular health-data generation, motivated by a documented sequence of null and inconsistent results from conventional rebalancing and filtering.
\item We introduce \sepa{}, a transparent post-generation selection rule that combines realism, label separability, and domain-anchor compatibility while preserving the real class distribution, so that any utility change cannot be attributed to a different class ratio.
\item We use a blocked $2^2$ factorial design under leakage-controlled five-fold cross-validation to attribute a synthetic-to-real utility change to separability pressure specifically, rather than to anchor pressure or to an unspecified combination of selection choices.
\item We report this effect across two health datasets with different sample sizes, imbalance ratios, and outcome heterogeneity, and show that it is strong and consistent for Vigitel obesity but statistically unsupported for COVID-19 mortality -- a dataset-dependent finding that we treat as informative rather than as a weakness to be hidden.
\end{enumerate}

The remainder of the paper is organised as follows. Section~\ref{sec:related-work} situates the work relative to tabular generative models, health-data evaluation, class-imbalance methods, and causally informed synthesis. Section~\ref{sec:methods} describes the datasets, leakage-controlled cross-validation protocol, the \sepa{} and \dacpf{} formulations, and the factorial experimental design. Section~\ref{sec:results} reports the empirical sequence, from the initial rebalancing and augmentation results through the fixed-$\lambda$ grid to the factorial attribution. Section~\ref{sec:discussion} interprets the proposed mechanism, considers the rival ``easier problem'' explanation, and discusses minority preservation, causal claims, privacy, and statistical limitations. Section~\ref{sec:limitations} converts the unresolved threats into a concrete set of confirmatory studies, and Section~\ref{sec:conclusion} states the present contribution and its boundaries.
